\documentclass[12 pt]{article}

\usepackage{stoversymb}
\usepackage[margin=0.75in, top=0.875in, bottom=0.875in]{geometry}
\usepackage{amsmath,amsfonts,amssymb,url,multicol,graphicx,tikz,soul,changepage}
%===makes urls render well===
\usepackage{lmodern}
\usepackage[T1]{fontenc}
%============================
\usepackage{wasysym} % smileys
\usepackage[inline]{enumitem}
\usepackage[many]{tcolorbox}
\tcbuselibrary{breakable}

%\everymath{\displaystyle}

\makeatletter
\renewcommand*\env@matrix[1][\arraystretch]{%
	\edef\arraystretch{#1}%
	\hskip -\arraycolsep
	\let\@ifnextchar\new@ifnextchar
	\array{*\c@MaxMatrixCols c}}
\makeatother

%\setenumerate{itemsep=0.25in}
%\setlist[enumerate,1]{label=\arabic*., itemsep=0.5in}
%\setlist[enumerate,2]{label={(\alph*)}, itemsep=0.8in}
%\setlist[enumerate,2]{leftmargin=0.5in}
%\setlist[enumerate,3]{label={\roman*.}}
\setlist[itemize,1]{label=$\circ$,topsep=0mm,itemsep=0mm}

\newcommand{\truefalse}[1]{#1\hfill\rule[-1mm]{220pt}{0.75pt}}
\newcommand{\hint}[1]{\vspace{3mm}\textbf{Hint}: #1}
\newcommand{\note}[1]{\textbf{Note}: #1}
\newcommand{\infsum}[3]{\sum_{{#1}={#2}}^\infty {#3}}
\newcommand{\axes}[1]{\begin{center}\includegraphics[scale=#1]{3DAxes}\end{center}}
\newcommand{\comps}[1]{\langle #1_1,#1_2,#1_3\rangle}
\newcommand{\compslong}[3]{\langle #1, #2, #3\rangle}
\newcommand{\ijk}[2]{#1\vect{#2}}
%\newcommand{\pmat}[1]{\begin{pmatrix}[1.25]#1\end{pmatrix}}
%\newcommand{\bmat}[1]{\begin{bmatrix}[1.25]#1\end{bmatrix}}
\newcounter{Exx}
\newcommand{\exbox}[1]{\stepcounter{Exx}\begin{tcolorbox}[breakable,arc=2pt,boxrule=1pt,right=6mm,top=3mm,bottom=6mm]\noindent\textit{\large\bfseries Example \theExx:\\[3mm]}~#1\end{tcolorbox}}
\newcommand{\exit}[1]{\vspace{4.5mm}\noindent\textit{Example #1:}}
\newcommand{\exbf}[1]{\vspace{1.5mm}\noindent\textbf{Example #1:}}
\newcommand{\LHS}{\text{LHS}}
\newcommand{\RHS}{\text{RHS}}
\newcommand{\ans}{\noindent\underline{Ans:}~}

\newenvironment{mypmatrix}[1]{\renewcommand{\arraystretch}{#1}\begin{pmatrix}}{\end{pmatrix}}
\newenvironment{mybmatrix}[1]{\renewcommand{\arraystretch}{#1}\begin{bmatrix}}{\end{bmatrix}}
\newcommand{\justify}{\ul{Justify your claim.}}

\newcommand{\pmat}[1]{\begin{mypmatrix}{1.25}#1\end{mypmatrix}}
\newcommand{\bmat}[1]{\begin{mybmatrix}{1.25}#1\end{mybmatrix}}

\makeatletter
\renewcommand*\env@matrix[1][*\c@MaxMatrixCols c]{%
	\hskip -\arraycolsep
	\let\@ifnextchar\new@ifnextchar
	\array{#1}}
\makeatother

\newcommand{\pmatgrid}[2]{\renewcommand{\arraystretch}{1.25}\begin{pmatrix}[#1] #2\end{pmatrix}}
\newcommand{\bmatgrid}[2]{\renewcommand{\arraystretch}{1.25}\begin{bmatrix}[#1] #2\end{bmatrix}}
\newcommand{\warnbox}[1]{\begin{adjustwidth}{0.375in}{0.375in}\begin{tcolorbox}[arc=0pt, colback=red!5!white, colframe=red, boxrule=0.5pt, parbox=false]{\parskip=3mm \parindent=0mm \red{\textbf{\large\itshape Warning!}\par #1}}\end{tcolorbox}\end{adjustwidth}}
\newcommand{\greentheorembox}[2]{\begin{adjustwidth}{0.375in}{0.375in}\begin{tcolorbox}[arc=0pt, colback=darkgreen!10!white, colframe=darkgreen, boxrule=0.5pt, parbox=false]{\parskip=3mm \parindent=0mm \green{\textbf{\large\itshape #1:}\par #2}}\end{tcolorbox}\end{adjustwidth}}
\newcounter{FactBox}
\newcommand{\factbox}[1]{\stepcounter{FactBox}\begin{adjustwidth}{0.375in}{0.375in}\begin{tcolorbox}[arc=0pt, colback=darkgreen!10!white, colframe=darkgreen, boxrule=0.5pt, parbox=false]{\parskip=3mm \parindent=0mm \green{\textbf{\large\itshape Fact \theFactBox:}\par #1}}\end{tcolorbox}\end{adjustwidth}}

\newcommand{\arrow}[1]{\xrightarrow{\hspace*{#1}}}

\setlength\parskip{0.75mm}	

\newcounter{ThmBox}
\newcommand{\thmbox}[1]{\stepcounter{ThmBox}\vspace{6mm}\noindent\fbox{\parbox{\linewidth}{~~(\theThmBox)~\hfill #1 \hfill\hfill}}}
\newcommand{\why}[1]{\begin{adjustwidth}{0.5in}{0.5in}{\par\vspace{3mm}\noindent\green{\textbf{Why should this be true? }#1}}\qed\end{adjustwidth}\vspace{6mm}}

\graphicspath{ {./../img/} }
\DeclareGraphicsExtensions{.pdf}

\begin{document}
%\begin{flushright}Name: \line(1,0){200}\end{flushright}
\begin{center}
\Large{\textbf{Properties of Transposes, Determinants, and Inverses}}
\end{center}
%\textbf{Directions:} Split into groups and (attempt to) answer each of the following questions with your group members. \ul{Don't worry about questions you are unable to answer!}

So far, we've learned how to take transposes, determinants, and inverses of (some) matrices; in this handout, we state some of the properties of these operations and provide some illustrative examples.

\subsection*{Transposes}
Recall that the \ul{transpose} of an $m\times n$ matrix $\sansA$ is the $n\times m$ matrix $\sansA^\sansT$ whose first row is the first column of $\sansA$, whose second row is the second column of $\sansA$, etc.

\exit{1} The transpose of $\pmat{1 & 2 & 3 \\ 4 & 5 & 6}$ is $\pmat{1 & 4 \\ 2 & 5 \\ 3 & 6}$.

\vspace{4.5mm}

\exit{2} The transpose of $\pmat{1 & 4 \\ 2 & 5 \\ 3 & 6}$ is $\pmat{1 & 2 & 3 \\ 4 & 5 & 6}$.

\vspace{4.5mm}

The above examples show one of the important properties of transposes (see T1 below) which is shown here with some others:

\vspace{4.5mm}

\noindent\fbox{\parbox{\linewidth}{
\begin{enumerate}[itemsep=0.25in,label=T\arabic*.,leftmargin=0.5in]
	\item $\left(\sansA^\sansT\right)^\sansT=\sansA$
	\item $\left(\sansA+\sansB\right)^\sansT=\sansA^\sansT+\sansB^\sansT$	
	\item $\left(\sansA\sansB\right)^\sansT=\sansB^\sansT\sansA^\sansT$ \hspace{3mm}\note{Notice the order swap on $\RHS$!}
\end{enumerate}
}}

\vspace{4.5mm}

\exit{3} If $\sansA=\pmat{1 & 2 \\ 3 & 4}$ and $\sansB=\pmat{5 & 6 \\ 7 & 8}$, then
\[
\LHS=\left(\sansA+\sansB\right)^\sansT=\left(\pmat{1 & 2 \\ 3 & 4}+\pmat{5 & 6 \\ 7 & 8}\right)^\sansT=\pmat{6 & 8 \\ 10 & 12}^\sansT=\pmat{6 & 10 \\ 8 & 12}
\]
while
\[
\RHS=\sansA^\sansT+\sansB^\sansT=\pmat{1 & 2 \\ 3 & 4}^\sansT+\pmat{5 & 6 \\ 7 & 8}^\sansT=\pmat{1 & 3 \\ 2 & 4}+\pmat{5 & 7 \\ 6 & 8}=\pmat{6 & 10 \\ 8 & 12},
\]

\noindent so $\LHS=\RHS$. \qed

\newpage

\exit{4} If $\sansA$ and $\sansB$ are as in Example 3, then $\LHS=\RHS=\pmat{19 & 43 \\ 22 & 50}$. \qed

\subsection*{Determinants}

Recall that $\det(\sansA)$ is a real number which is defined if and only if $\sansA$ is a square matrix and that---to compute $\det(\sansA)$---it suffices to do ``cofactor expansion'' along any row or column of $\sansA$. See the lecture notes if you need a refresher on how to do this!

There are two main properties of determinants that you should know.

\vspace{4.5mm}

\noindent\fbox{\parbox{\linewidth}{
\begin{enumerate}[itemsep=0.25in,label=D\arabic*.,leftmargin=0.5in]
	\item $\det(\sansA^\sansT)=\det(\sansA)$
	\item $\det(\sansA\sansB)=\det(\sansA)\det(\sansB)$
\end{enumerate}
}}

\vspace{4.5mm}

\exit{5} If $\sansA$ is as in Example 3, then 
\[
\LHS=\det\pmat{1 & 2 \\ 3 & 4}^\sansT=\det\pmat{1 & 3 \\ 2 & 4}=4-6=-2
\]
while
\[
\RHS=\det\pmat{1 & 2 \\ 3 & 4}=4-6=-2,
\]

\noindent so $\LHS=\RHS$. \qed

\vspace{9mm}

\exit{6} If $\sansA=\pmat{1 & 2 \\ 3 & 4}$ and $\sansB=\pmat{5 & 6 \\ 7 & 8}$ as in Example 3, then $\sansA\sansB=\pmat{19 & 22 \\ 43 & 50}$ and so
\[
\LHS=\det(\sansA\sansB)=\det\pmat{19 & 22 \\ 43 & 50}=19(50)=43(22)=4
\]
while
\[
\RHS=\det(\sansA)\det(\sansB)=\det\pmat{1 & 2 \\ 3 & 4}\cdot\det\pmat{5 & 6 \\ 7 & 8}=-2\cdot-2=4.
\]

\noindent Thus, $\LHS=\RHS$. \qed

\newpage

\subsection*{Inverses}

With determinants in tow, we're able to define and talk about matrix inverses!

Recall that the \ul{inverse} of an $n\times n$ matrix $\sansA$ is an $n\times n$ matrix $\sansA^{-1}$ for which $\sansA\sansA^{-1}=\sansI_n=\sansA^{-1}\sansA$ where $\sansI_n$ is the $n\times n$ identity matrix. 

As we learned in class, $\sansA$ has an inverse if and only if $\det(\sansA)\neq 0$ and---when it exists---you can find $\sansA^{-1}$ by forming the augmented matrix $\pmatgrid{c|c}{\sansA & \sansI_n}$ and putting the ``$\sansA$ part'' into RREF. The result is guaranteed to be the augmented matrix 
\[
	\pmatgrid{c|c}{\sansI_n & \sansA^{-1}}.
\]

In a future handout, we'll learn a whole lot more about matrix inverses; until then, here are some properties you should know!

\vspace{4.5mm}

\noindent\fbox{\parbox{\linewidth}{
\begin{enumerate}[itemsep=0.25in,label=I\arabic*.,leftmargin=0.5in]
	\item $(\sansA^{-1})^{-1}=\sansA$
	\item $\left(\sansA\sansB\right)^{-1}=\sansB^{-1}\sansA^{-1}$.\hspace{3mm}\note{Notice the order swap on $\RHS$!}
	\item $\left(\sansA^\sansT\right)^{-1}=\left(\sansA^{-1}\right)^\sansT$.
\end{enumerate}
}}

\vspace{4.5mm}

\exit{7} If $\sansA$ is as in Example 3, then $\sansA^{-1}=\pmat{-2 & 1 \\ \frac{3}{2} & -\frac{1}{2}}$ and the inverse of $\sansA^{-1}$ is precisely $\sansA$. \textbf{You should verify this!} \qed

\vspace{9mm}

\exit{8} If $\sansA=\pmat{1 & 2 \\ 3 & 4}$ and $\sansB=\pmat{5 & 6 \\ 7 & 8}$ as above, then $\sansA\sansB=\pmat{19 & 22 \\ 43 & 50}$,  \mbox{$\sansA^{-1}=\pmat{-2 & 1 \\ \frac{3}{2} & -\frac{1}{2}}$}, and $\sansB^{-1}=\pmat{-4 & 3 \\ \frac{7}{2} & -\frac{5}{2}}$. Thus,
\[
\LHS=(\sansA\sansB)^{-1}=\det\pmat{19 & 22 \\ 43 & 50}^{-1}=\pmat{\frac{25}{2} & -\frac{11}{2} \\ -\frac{43}{4} & \frac{19}{4}}
\]
\noindent while
\[
\RHS=\sansB^{-1}\sansA^{-1}=\pmat{-4 & 3 \\ \frac{7}{2} & -\frac{5}{2}}\pmat{-2 & 1 \\ \frac{3}{2} & -\frac{1}{2}}=\pmat{\frac{25}{2} & -\frac{11}{2} \\ -\frac{43}{4} & \frac{19}{4}},
\]

\noindent and so $\LHS=\RHS$. \textbf{You should verify this, too!}\qed

\vspace{9mm}

\exit{9} If $\sansA$ is as in Example 3, then $\sansA^\sansT=\pmat{1 & 3 \\ 2 & 4}$ has inverse $\LHS=\pmat{-2 & \frac{3}{2} \\ 1 & -\frac{1}{2}}$. On the other hand, $\sansA^{-1}=\pmat{-2 & 1 \\ \frac{3}{2} & -\frac{1}{2}}$ (see Example 7) and so 
\[
\RHS=\left(\sansA^{-1}\right)^\sansT=\pmat{-2 & 1 \\ \frac{3}{2} & -\frac{1}{2}}^\sansT=\pmat{-2 & \frac{3}{2} \\ 1 & -\frac{1}{2}}.
\]
\noindent Thus, $\LHS=\RHS$. \textbf{Verify everything!}\qed
\end{document}